% Einheit 3 von 3 – Skalarprodukt im Sachzusammenhang (bank/vektoren-und-rechenoperationen/e3.jsonl, zusammenbau v0.1)
%% TODO Zeitmarke Einheit 3: Marken-Zeile nicht lesbar
%% TODO Zweigzeile Teil 1: was hier gelernt wird (Satz in Schülersprache aus dem Einheitstitel) – Einheit 3
\einheitenkopf[e3]{Einheit 3 von 3 · Skalarprodukt im Sachzusammenhang}
\zweigzeile{keine Prüfungsaufgabe}
\setcounter{aufgabe}{13}

%% TODO Titel: Ich-kann-Satz fehlt in der Bank (Platzhalter: Kettenname) – Nr. 14
%% TODO Anweisung: ein Satz über der ersten Teilaufgabe fehlt in der Bank – Nr. 14
\begin{aufgabe}{Skalarprodukt im Sachzusammenhang}
%% TODO \rechenplatz{n}: Schrittzahl des Grundfalls fehlt in der Bank (nur bei fester Schreibform, 2.3 b)
\begin{teile}
\teil Die Liste der gekauften Mengen je Sorte (Äpfel, Birnen, Pflaumen): Zahl oder Vektor? \\ \kreuz{Zahl} \\ \kreuz{Vektor}
\teil Ein Café verkauft an einem Tag 12 Espresso, 30 Cappuccino und 18 Tee. Absatzvektor (Espresso | Cappuccino | Tee)? Was zählt die zweite Komponente? ( \leerfeld | \leerfeld | \leerfeld )
\teil Ein Imbiss kauft $\vec m = (k | s | w)$ Kisten Cola, Saft und Wasser; $\vec p = (14{,}50 | 18{,}00 | 6{,}40)$ sind die Preise in Euro je Kiste. Was gibt $\vec m \circ \vec p$ an? \hfill (Abitur 2019 GK)
\teil Ein Maler mischt Weiß, Gelb und Rot im Verhältnis 5 : 2 : 1; $\vec m = (w | g | r)$ in Litern, $\vec p = (4 | 6 | 8)$ Euro je Liter, $\vec m \circ \vec p = 168$. Wie viel Liter jeder Farbe? \hfill (Abitur 2019 GK)
\end{teile}
\end{aufgabe}

%% TODO Titel: Ich-kann-Satz fehlt in der Bank (Platzhalter: Kettenname) – Nr. 15
%% TODO Anweisung: ein Satz über der ersten Teilaufgabe fehlt in der Bank – Nr. 15
\begin{aufgabe}{Skalarprodukt im Sachzusammenhang – Fehler finden, Begründen, Anwendung}
\begin{teile}
\teil Ich finde den Fehler: $\vec m = (2 | 3 | 1)$ kg, $\vec p = (4 | 2 | 7)$ Euro je kg. Nora: $\vec m \circ \vec p = (8 | 6 | 7)$ €.
\teil Warum ist $\vec m \circ \vec p$ eine Zahl und kein Vektor?
\teil Eine Klasse bestellt 24 Hefte zu 1,20 €, 12 Blöcke zu 2,50 € und 30 Stifte zu 0,80 €. In der Klassenkasse sind 80 €. Reicht das?
\end{teile}
\end{aufgabe}
